\documentclass[aspectratio=169]{beamer}
\usetheme{Madrid}
\usepackage{amsmath, amssymb}

\title{Functions -- Finally Made Simple}
\subtitle{CSCE 222: Discrete Structures for Computing}
\author{}
\date{}

\begin{document}

%------------------------------------------------
\begin{frame}
\titlepage
\end{frame}

%------------------------------------------------
\begin{frame}{First: Breathe. This Is Not Hard.}
You are not confused because this is hard.  

You are confused because it was presented without structure.

Today we:
\begin{itemize}
\item Slow everything down.
\item Define every symbol the first time we use it.
\item Show why every condition exists.
\item Connect everything together.
\end{itemize}

Functions are just:
\[
\textbf{Disciplined input-output rules.}
\]

That's it.
\end{frame}

%------------------------------------------------
\begin{frame}{Cheat Sheet (Part 1): Core Definition}
\begin{columns}
\column{0.5\textwidth}
\textbf{Function Definition}

A function $f$ from set $A$ to set $B$ is:

\[
f : A \to B
\]

Such that:

\[
\forall a \in A, \exists! \, b \in B
\text{ where } f(a) = b
\]

\column{0.5\textwidth}
\textbf{What Every Symbol Means}

$A$ = domain (inputs)

$B$ = codomain (allowed outputs)

$a$ = element of $A$

$b$ = element of $B$

$\exists!$ = \emph{exactly one}

If ``exactly one'' did not exist:

One input could map to two outputs.

Then it would not be predictable.

Not a function.
\end{columns}
\end{frame}

%------------------------------------------------
\begin{frame}{Cheat Sheet (Part 2): Terminology}
\begin{columns}
\column{0.5\textwidth}
For $f : A \to B$ and $f(a)=b$

\textbf{Domain}: $A$

\textbf{Codomain}: $B$

\textbf{Image of $a$}: $b$

\textbf{Preimage of $b$}: all $a$ such that $f(a)=b$

\column{0.5\textwidth}
\textbf{Range}

Range = all actual outputs.

\[
\text{Range} = \{ f(a) \mid a \in A \}
\]

Codomain = allowed outputs.

Range = actual outputs.

They are NOT always equal.
\end{columns}
\end{frame}

%------------------------------------------------
\begin{frame}{Cheat Sheet (Part 3): Function Equality}
\begin{columns}
\column{0.5\textwidth}
Given:

\[
F: X \to Y
\]
\[
G: X \to Y
\]

\[
F = G \iff F(x)=G(x)
\quad \forall x \in X
\]

\column{0.5\textwidth}
Meaning:

Same domain.

Same codomain.

Same output for every input.

If even ONE input differs:

Functions are not equal.
\end{columns}
\end{frame}

%------------------------------------------------
\begin{frame}{Cheat Sheet (Part 4): One-to-One}
\begin{columns}
\column{0.5\textwidth}
$F$ is one-to-one (injective) iff:

\[
F(x_1)=F(x_2) \Rightarrow x_1=x_2
\]

Equivalent:

\[
x_1 \ne x_2 \Rightarrow F(x_1)\ne F(x_2)
\]

\column{0.5\textwidth}
Why this matters:

Different inputs give different outputs.

If two inputs share output:

Not injective.

Think: No collapsing.
\end{columns}
\end{frame}

%------------------------------------------------
\begin{frame}{Cheat Sheet (Part 5): Onto}
\begin{columns}
\column{0.5\textwidth}
$F$ is onto (surjective) iff:

\[
\forall y \in Y,
\exists x \in X
\text{ such that } F(x)=y
\]

\column{0.5\textwidth}
Meaning:

Every element in codomain is used.

If one $y$ never appears:

Not onto.

Think: No unused outputs.
\end{columns}
\end{frame}

%------------------------------------------------
\begin{frame}{Cheat Sheet (Part 6): Bijection and Inverse}
\begin{columns}
\column{0.5\textwidth}
Bijection:

One-to-one AND onto.

\[
F: X \to Y
\]

If bijective, then inverse exists:

\[
F^{-1}: Y \to X
\]

\column{0.5\textwidth}
Inverse condition:

\[
F^{-1}(F(x)) = x
\]

Inverse only exists if:

No collapsing AND no unused outputs.
\end{columns}
\end{frame}

%------------------------------------------------
\begin{frame}{Cheat Sheet (Part 7): Floor and Ceiling}
\begin{columns}
\column{0.5\textwidth}
For real $x$:

\[
\lfloor x \rfloor
\]

Largest integer $\le x$

\[
\lceil x \rceil
\]

Smallest integer $\ge x$

\column{0.5\textwidth}
Example:

$x = 3.2$

\[
\lfloor 3.2 \rfloor = 3
\]
\[
\lceil 3.2 \rceil = 4
\]

Important for counting objects.
\end{columns}
\end{frame}

%------------------------------------------------
\begin{frame}{Now Let’s Build Intuition}
A function is just:

\[
\text{Machine: input } \to \text{ output}
\]

The rule:

\[
f(x) = x^2
\]

Means:

Take input $x$.

Multiply it by itself.

Return result.
\end{frame}

%------------------------------------------------
\begin{frame}{Example: Is This a Function?}
Let:

\[
A = \{3,5,7\}
\]
\[
B = \{2,4\}
\]

Relation:

\[
R = \{(3,2),(3,4),(7,2)\}
\]

Why not a function?

Input 3 maps to TWO outputs.

Violates ``exactly one''.
\end{frame}

%------------------------------------------------
\begin{frame}{One-to-One Example}
Let:

\[
f(x)=x^2, \quad f:\mathbb{Z}\to\mathbb{Z}
\]

Test injective:

\[
f(4)=16
\]
\[
f(-4)=16
\]

Two different inputs.

Same output.

Therefore:

Not one-to-one.
\end{frame}

%------------------------------------------------
\begin{frame}{Onto Example}
Same function:

\[
f(x)=x^2, \quad f:\mathbb{Z}\to\mathbb{Z}
\]

Is it onto?

Is there $x$ such that:

\[
f(x)=2?
\]

No integer squared gives 2.

So not onto.
\end{frame}

%------------------------------------------------
\begin{frame}{Inverse Example (Step 1)}
Let:

\[
F(x)=4x-5
\]

Step 1: Replace $F(x)$ with $y$:

\[
y=4x-5
\]

Goal: Solve for $x$.
\end{frame}

%------------------------------------------------
\begin{frame}{Inverse Example (Step 2)}
Solve:

\[
y=4x-5
\]

Add 5:

\[
y+5=4x
\]

Divide by 4:

\[
x=\frac{y+5}{4}
\]

So:

\[
F^{-1}(y)=\frac{y+5}{4}
\]
\end{frame}

%------------------------------------------------
\begin{frame}{Why Floor and Ceiling Matter}
Example:

100 bits.

1 byte = 8 bits.

\[
\frac{100}{8}=12.5
\]

You cannot store half a byte.

So:

\[
\lceil 12.5 \rceil = 13
\]

Ceiling because we must round up.
\end{frame}

%------------------------------------------------
\begin{frame}{Homework Survival Strategy}
Every problem asks one of four things:

\begin{itemize}
\item Is it a function?
\item Is it injective?
\item Is it onto?
\item Does inverse exist?
\end{itemize}

Translate the problem into definitions.

Then apply logic slowly.

No tricks. Just definitions.
\end{frame}

%------------------------------------------------
\begin{frame}{Extra Example (Deep Why)}
Let:

\[
f: \mathbb{R} \to \mathbb{R}
\]
\[
f(x)=3x+2
\]

Injective?

Assume:

\[
f(x_1)=f(x_2)
\]

\[
3x_1+2=3x_2+2
\]

Subtract 2:

\[
3x_1=3x_2
\]

Divide 3:

\[
x_1=x_2
\]

Therefore injective.
\end{frame}

%------------------------------------------------
\begin{frame}{Final Perspective}
Functions are:

Structured mappings.

Injective = no collapsing.

Onto = no unused outputs.

Bijection = perfect matching.

Inverse = undo button.

This is pattern recognition.

Not chaos.

You absolutely can do this.
\end{frame}

\end{document}